\appendix

\revised{
\section{\sysname-100 Benchmark}
\label{app:benchmark}
\subsection{Benchmark overview}
\label{app:benchmark_overview}
\begin{figure}[h]
    \centering
    \resizebox{0.95\linewidth}{!}{\includegraphics{Figures/ava100-overview.pdf}}
    \caption{\revised{Scenarios and QA examples in \sysname-100.}}
    \label{fig:ava-100-overview}
\end{figure}

\sysname-100 is established as a benchmark to advance research in L4 video analytics systems, with a particular focus on ultra-long-duration video understanding and complex reasoning tasks. In contrast to existing datasets that primarily focus on short video clips or domain-specific activities, \sysname-100 centers on real-world, continuous video streams often spanning several hours. This design enables the evaluation of video analytics systems under realistic and challenging scenarios, including long-horizon temporal reasoning, sparse event detection, and cross-segment summarization, etc.

The benchmark currently comprises eight ultra-long videos encompassing diverse scenarios including egocentric recordings of daily activities, city tours, as well as third-person fixed-camera capturing traffic monitoring and wildlife monitoring. Each video is annotated with multiple question–answer (QA) pairs: human annotators generate ground-truth QAs, while GPT-4o is employed to produce multiple-choice options, subsequently verified through manual review. In total, \sysname-100 contains approximately 100 hours of video content and 120 annotated QA pairs, offering a valuable resource for the evaluation of both video-language models and specialized video analytics systems.

Fig.~\ref{fig:ava-100-overview} illustrates the representative scenarios and QA examples featured in \sysname-100. Table~\ref{tab:ava100} presents detailed statistics of the \sysname-100 dataset, encompassing video duration and the number of QA pairs. The benchmark is publicly available at \url{https://huggingface.co/datasets/iesc/Ava-100}.
}


\revised{
\subsection{Selected Scenarios and Data Sources}
\subsubsection{Human Daily Activities}


The egocentric daily-life scenario videos in \sysname-100 are primarily sourced from the Ego4D dataset~\cite{grauman2022ego4d}. Ego4D is one of the largest existing egocentric video datasets, containing more than 3,600 hours of video collected across multiple countries, focusing on a wide range of daily activities such as household routines, cooking, shopping, and social interactions. It was originally designed to support research in egocentric perception, including tasks like action recognition, episodic memory, forecasting, and hand–object interaction understanding.  

\begin{table}[h]
\centering
\resizebox{\linewidth}{!}{
\begin{tabular}{lccc}
\toprule[1.5pt]
\textbf{Video ID} & \textbf{Duration (hours)} & \textbf{\#QA Pairs} & \textbf{Views} \\
\midrule
ego-1 & 12.7 & 22 & First-person (moving) \\
ego-2 & 11.7 & 19 & First-person (moving) \\
citytour-1 & 12.0 & 19 & First-person (moving) \\
citytour-2 & 10.5 & 20 & First-person (moving) \\
\midrule
traffic-1 & 14.9 & 12 & Third-person (fixed) \\
traffic-2 & 13.9 & 13 & Third-person (fixed) \\
wildlife-1 & 12.0 & 8  & Third-person (fixed) \\
wildlife-2 & 11.5 & 7  & Third-person (fixed) \\
\midrule
\textbf{Total} & \textbf{99.2} & \textbf{120} & -- \\
\bottomrule[1.5pt]
\end{tabular}
}
\caption{\revised{Statistics of the \sysname-100 dataset, including video duration, number of QA pairs, and perspective attributes for each video. The first four entries correspond to egocentric videos capturing moving perspectives, whereas the last four represent fixed third-person recordings.}}
\label{tab:ava100}
\end{table}

Ego4D does not typically provide single, continuous clips exceeding 10 hours in duration. To construct long-form egocentric videos for \sysname-100, we employed the concatenation strategy. Specifically, we prioritized the preservation of event diversity over maintaining the identity of the camera wearer. Accordingly, we selected and concatenated sub-clips from distinct Ego4D videos to generate extended sequences surpassing 10 hours. Throughout this process, we guarantee that the concatenated sub-clips encompassed a wide spectrum of human activities and contexts, thereby enhancing the richness of the video content while mitigating potential annotation ambiguities.
}

\revised{
\subsubsection{City Walking}

The first-person city tour videos in \sysname-100 are sourced from the YouTube channel \textit{4K World Wandering}\cite{4KWorldWandering}, which features high-resolution walking tour videos recorded in various locations around the globe.  These videos capture natural human navigation experiences in public spaces, making them a valuable complement to the daily-life egocentric content. The videos typically span up to four hours in length. To construct long-form sequences for \sysname-100, we adopted the same concatenation strategy as described above. Specifically, we aimed to compose extended videos by combining multiple sub-clips, while preserving both geographic and contextual coherence.  

In particular, we curated sub-videos predominantly from the same region or country to construct each long-form video. Meanwhile we intentionally enhanced diversity by integrating segments filmed under various weather conditions (\eg sunny, snowy, rainy) and at different times of day (\eg daytime, night) into the concatenated sequences. This design enriches the video contents, rendering the city walking videos in \sysname-100 more representative of real-world navigation scenarios across diverse geographic and temporal conditions.


The city walking videos prioritize large-scale spatial cognition, long-term trajectory tracking, and environmental variability. Consequently, they present added challenges for QA annotation, as the questions frequently require reasoning about landmarks, navigation routes, or temporal dynamics within the environment.
}

\revised{
\subsubsection{Traffic Monitoring}

The traffic monitoring videos in \sysname-100 are sourced from the Bellevue Traffic dataset~\cite{bhardwaj2022ekya}, which comprises continuous, fixed-camera recordings of real-world traffic scenes. This dataset is particularly well-suited for the study of long-term temporal dynamics, vehicle-pedestrian interactions, and congestion analysis. 

For \sysname-100, we selected two intersections: Bellevue\_150th\_Newport (spanning from 2017-09-11 03:08:29 to 2017-09-11 16:08:32) and Bellevue\_Bellevue\_NE8th (spanning from 2017-09-11 03:08:29 to 2017-09-11 15:08:32). These extensive, continuous recordings encompass a broad spectrum of traffic events across different periods of the day, including both peak hours and quieter intervals, thereby providing a rich and diverse dataset.

To construct QA pairs, we target the fine-grained traffic events, such as vehicle movements, pedestrian crossings, and congestion incidents. For example, whether a red car passed through the intersection between 4:30 and 4:40 and in which direction it traveled; the presence of pedestrians or cyclists at specific times. Such temporally anchored, detail-oriented QAs highlight the unique value of traffic monitoring videos for evaluating long-duration reasoning.

}

\revised{
\subsubsection{Wildlife Monitoring}

The wildlife monitoring videos in \sysname-100 are sourced from publicly available YouTube live-streaming channels, which continuously capture activities in natural habitats through fixed outdoor cameras. Particularly, we select four sources. The first source is the Arkansas Critter Cam Channel\cite{ArkansasCritterCam}, from which we used the video titled "Wildlife Live – African Waterhole 24/7", originally published on June 17, 2020. This stream features footage from a wildlife camera located in Arkansas, USA.  The second source is the "Nature Live Cams" channel\cite{NatureLive}, which specializes in streaming diverse natural environments worldwide, including forests, rivers, and savannahs. We selected the video "African Safari Live 24/7 – Watering Hole and Savannah" and extracted two segments from its April 12, 2025 livestream: 05:03–11:03 and 17:49–23:49. These segments were concatenated to form one ultra-long video. We try to ensure the coverage of animal activities across different times of the day, from early morning to late night. It is worth noting that wildlife scenes are inherently random and unpredictable, resulting in key events occurring infrequently and at indeterminate times within long-duration videos, posing additional challenges for L4 video analytics systems.
}

\revised{
\subsection{Prompts}
\label{app:prompts}


In \sysname-100, we employ a scenario-specific prompt design strategy, predicated on the view that prompts should be tailored to the deployment contexts of video analytics systems in order to optimize VLM performance by extracting salient information under varying conditions. For instance, prompts for first-person daily life videos are tailored to emphasize object interactions and activity sequences, whereas prompts for traffic monitoring footage prioritize vehicle trajectories, traffic flow, and pedestrian dynamics. Crucially, we regard prompt optimization as an integral component of system-level optimization within a fully VLM-driven video analytics framework. Thoughtfully crafted prompts can substantially mitigate redundant reasoning, enhance the efficiency of long-range temporal information extraction, and thereby elevate overall system performance. 

The prompts employed in our experiments are detailed as follows.





}

\balance

\clearpage
\begin{figure*}[t]
\centering
\begin{minipage}{0.9\textwidth}
\begin{lstlisting}[language=, caption={Description generation prompt for Human Daily Activities Scenario}, basicstyle=\footnotesize]
You are an expert in video understanding and description generation. 
You are given a first-person perspective video, and your task is to generate a continuous, smooth, and grounded description of the video content.

Focus particularly on:
- The actions and events performed by the camera wearer (the person holding or wearing the camera).
- The surrounding environment, including objects, people, and notable visual changes.
- The physical characteristics and spatial relationships of objects in the environment (e.g., size, color, relative positions, proximity to the camera wearer).
- Interactions between the camera wearer and the environment, including object manipulations and movements through space.

Avoid describing each frame individually, such as "frame1...". Instead, provide a coherent and logically structured narrative that flows smoothly over time.

Important constraints:
- Do not include assumptions, inferences, or fabricated details that are not visually evident.
- Do not speculate about the identity, emotions, or intentions of the camera wearer unless explicitly shown.
- When referring to the person holding or wearing the camera, always use the term "camera wearer".

Return your response as a single, continuous, and fluent paragraph that comprehensively describes the video content, including fine-grained visual details, and limit the length to 400 words.
\end{lstlisting}
\end{minipage}
\end{figure*}

\begin{figure*}[t]
\centering
\begin{minipage}{0.9\textwidth}
\begin{lstlisting}[language=, caption={Description generation prompt for City Walking Scenario}, basicstyle=\footnotesize]
You are an expert in video understanding and detailed scene description. 
You are given a first-person perspective video of a person walking through a city environment and your task is to generate a continuous, smooth, and grounded description of the video content.

Focus particularly on:
- The locations and landmarks the camera wearer passes by, such as buildings, shops, streets, and public spaces.
- The appearance and functions of these places (e.g., a small bakery with a red awning, a tall glass office building, a busy intersection).
- Events or notable occurrences observed during the walk, such as street performances, traffic changes, or people interacting in public.

Avoid describing each frame individually. Instead, provide a logically structured narrative that flows naturally over time.

Important constraints:
- Do not include assumptions, inferences, or fabricated details that are not visually evident.
- Do not speculate about the identity, emotions, or intentions of the camera wearer or other people unless explicitly shown.
- When referring to the person holding or wearing the camera, always use the term "camera wearer".

Return your response as a single, continuous, and fluent paragraph that comprehensively describes the video content, with attention to fine-grained urban and visual details, and limit the length to 400 words.
\end{lstlisting}
\end{minipage}
\end{figure*}

\begin{figure*}[t]
\centering
\begin{minipage}{0.9\textwidth}
\begin{lstlisting}[language=, caption={Description generation prompt for Traffic Monitoring Scenario}, basicstyle=\footnotesize]
You are a video analysis expert specializing in traffic observation and detailed event description. You are analyzing a road or intersection monitoring video recorded by a fixed-position camera.
Your task is to generate an **accurate, grounded, and coherent** description of the video segment.
Please focus on the following aspects:
- **Observed Traffic Elements:** Identify all traffic-related elements present in the video. Provide an integrated description covering the entire segment:
- **Vehicle Types:** Identify types as accurately as possible (e.g., car, truck, bus, motorcycle, bicycle, van). If unclear, describe the vehicle's physical characteristics (e.g., "a large box truck," "a small passenger vehicle," "a two-wheeled vehicle").
- **Quantity:** Indicate the number of each identified vehicle type, as well as the number of pedestrians.
- **Characteristics:** If relevant to the scene, note distinguishing physical features (e.g., color, size, presence of trailers, specific structural features). Describe pedestrians based on their interaction with traffic (e.g., walking along the sidewalk, crossing the street).
- **Actions / Events:** Describe observed dynamic behaviors and interactions (e.g., driving in a specific lane, stopping, turning, entering/exiting the frame, changing lanes, overtaking, pedestrians waiting or crossing), including any **traffic anomalies** (e.g., sudden braking, erratic maneuvers, red-light violations, collisions, traffic violations, illegal parking that obstructs traffic).
- **Timestamps:** Identify the timestamp shown on the monitoring footage.
**Output Format:**
After watching the full video segment, write a structured summary paragraph in the following format:
[Timestamp]: [Summary of vehicle types, quantities, characteristics, actions, pedestrian activity, and traffic anomalies].
**Important Constraints:**
- Do not include assumptions or details not clearly visible in the footage.
- Do not speculate about the intentions or emotions of drivers or pedestrians; only describe observable actions and postures.
- Use neutral descriptions such as "the footage shows" or "the camera captures"; avoid subjective phrasing like "we see" or "the viewer can observe."
- If vehicle type identification is uncertain, state this clearly.
- The final output should be a concise, fact-based summary of the traffic activity and scene context. The length should be appropriate to the events observed with a target length of less than 400 words.
\end{lstlisting}
\end{minipage}
\end{figure*}

\begin{figure*}[t]
\centering
\begin{minipage}{0.9\textwidth}
\begin{lstlisting}[language=, caption={Description generation prompt for Wildlife monitoring Scenario}, basicstyle=\footnotesize]
You are an expert in video analysis, specializing in wildlife observation and detailed environmental description.  
You are analyzing fixed-camera monitoring footage capturing a scene in a wild or natural environment and your task is to generate a precise, grounded, and chronologically ordered description of the entire video segment.

Focus on the following aspects:
- **Observed Animals:** Identify any animals present in the footage. For the entire segment, provide a consolidated description of:
- **Species:** Identify species as accurately as possible. If uncertain, describe physical characteristics (e.g., "a large brown bear", "a small rodent-like mammal", "a flock of unidentified birds").
- **Number:** Indicate the number of individuals observed.
- **Appearance:** Note distinctive physical features (e.g., size, color, antlers, markings).
- **Behavior:** Describe observed behaviors (e.g., foraging, running, resting, entering/exiting the frame, interacting).
- **Timestamps:** Identify the timestamp displayed in the monitoring footage.
- **Environment:** Briefly describe the static environment visible to the camera (e.g., forest clearing, rocky terrain, vegetation), and note any significant changes during the observation period (e.g., lighting shifts, weather changes).

**Output Format:**  
After reviewing the full segment, summarize your findings in a single structured paragraph using the following format:

[Timestamp]: [Environment description][Summary of animal and their activities]

**Important Constraints:**
- Do not include assumptions or invented details that are not visually evident in the footage.
- Do not speculate on the intentions or emotions of the animals; describe only observable actions and postures.
- Refer to observations using neutral terms such as "the footage shows" or "the camera captures"; avoid subjective phrasing like "we see" or "the viewer can observe".
- If species identification is uncertain, explicitly state this.
- The final output should be a concise, fact-based summary of the wildlife activity and environmental context of the segment, with a target length of less than 400 words.
\end{lstlisting}
\end{minipage}
\end{figure*}

